\documentclass[options]{philosophersimprint}
%\usepackage{sectsty}
%\linespread{1.2}
\usepackage{xunicode}
\usepackage{xltxtra}
\defaultfontfeatures{Mapping=tex-text}
\setmainfont[
Ligatures={Common}, Numbers={OldStyle}, Variant=01,
BoldFont=LinLibertine_RB.otf,
ItalicFont=LinLibertine_RI.otf,
BoldItalicFont=LinLibertine_RBI.otf
]{LinLibertine_R.otf}
\setmonofont[Scale=0.8]{DejaVuSansMono.ttf}
\usepackage{stmaryrd}
\usepackage[toc,page]{appendix}
\usepackage{amsmath,amsthm,amssymb}
\usepackage[colorlinks=true, citecolor=blue, linkcolor=blue, hyperfootnotes=false]{hyperref}
\usepackage{nameref}
\usepackage{soul}
\usepackage{frame}
\usepackage[style=apa, natbib=true]{biblatex}
\usepackage{scrextend}
\usepackage{mathrsfs}
%\addtokomafont{labelinglabel}{\sffamily}
\makeatletter
\newcommand\labelitem[2]{%
  \item[#1]\def\@currentlabel{#1}\label{#2}}
\makeatother
\usepackage[inline]{enumitem}
\usepackage{framed, csquotes}
\usepackage{tikz}
\usepackage{scalerel}
\usepackage{titlesec}
\newcommand\halo{{{\circ}\mkern-1mu}}
\usepackage{linguex}
\usepackage{parskip}
\usepackage{fancyhdr}
\pagestyle{plain}
% Diagrams
%\usepackage{pgfplots}
\usepackage{tikz}
\usepackage{qtree}




\usepackage{FOL_Yale/notation}

\begin{document}

PHIL 1115: First-order Logic \hfill Yale, Fall 2026\\
Lecture 5, Handout \hfill Dr. Daniel Grimmer

\textbf{1. Taking Stock} So far in this course, we've been thinking about notions like truth and falsity, validity and invalidity, and so on. To help get a handle on these notions, we've built an artificial, formal language with which to study them: the language of propositional logic.

We've also come up with a systematic procedure for telling us when a complex formula of propositional logic is true/false as well as when an argument is valid. The procedure in question here involves making \textit{truth tables}. In a sense, then, what we've done is built up a \textbf{logical system}. We'll be adding to this system in the second half of the course. This more general system will be called \textit{first-order logic}. In subsequent logic courses you could extend this system further: e.g., to modal logic or second-order logic.

But in the meantime, one of the things we can do now is start asking questions \textit{about} our system. That is, we can begin doing some \textit{meta-logic}. In the real world, whenever a logician comes up with a system like ours, they immediately start asking questions of this kind. For example: What is the scope of the system? What are its limits? What are its invariant structures?

\textbf{Example Question}. Consider some generic formulas. For each formula, count its atomic formulas (including repetitions) as well as the number of \textit{binary} connectives (\textit{binary} here means $\ConjTight$, $\Disj$, $\Cond$, and $\Bicond$ but not $\Neg$).
\begin{itemize}
    \item \makebox[3.7cm][l]{$\Neg p$} 1 atomic, 0 binary connectives
    \item \makebox[3.7cm][l]{$(p \ \Bicond \ \Neg q)$} 2 atomics, 1 binary connective
    \item \makebox[3.7cm][l]{$((p \ \Cond \ q)\Disj s)$} 3 atomics, 2 binary connectives
    \item \makebox[3.7cm][l]{$((p \Conj q)\Cond(p \Disj s))$} 4 atomics, 3 binary connectives
    \item \makebox[3.7cm][l]{$(s\Conj(((p\Cond q)\Cond r)\Bicond q))$} 5 atomics, 4 binary connectives
\end{itemize}
\textit{Conjecture}: the number of atomic formulas is always one greater than the number of binary connectives.

\textbf{Challenge Question}. Consider the truth tables for the formulas $(p \Conj q)$, $(p \Disj q)$, $(p\Cond q)$, and $(p \Bicond q)$. You'll notice here that each of these formulas is true in the \textbf{all-true model}: the first row of its truth table, where every atomic formula is true. \textit{Conjecture}: If a formula doesn't contain $\Neg$ then it is \textit{always} true in the all-true model.

These two conjectures seem plausible. But how can we know \textit{for sure} whether they are true? (We prove them below, as Examples 3 and 4.)

\textbf{2. Mathematical Induction!} One of the most powerful techniques we can use to prove things \textit{about} our logical system is the technique known as \textbf{mathematical induction}.\footnote{A warning about the word. \textit{Mathematical} induction has nothing to do with \textit{inductive} reasoning (the risky, probabilistic kind Restall sets aside in Chapter 1, and the subject of our Lecture 23). A proof by mathematical induction is entirely deductive: airtight, not merely probable. The two share a name and nothing else.} Here are two examples (one kind of boring, and the other more fun):

\textbf{Example 1}. Suppose you have infinitely many dominoes lined up. An infinite number of conditional claims are then true:
\begin{itemize}
    \item if domino $1$ falls over, then so does domino $2$
    \item if domino $2$ falls over, then so does domino $3$
    \item if domino $3$ falls over, then so does domino $4$
    \item \dots \ (forever)
\end{itemize}
In general: if an \textit{arbitrary} domino $n$ falls over, then so does domino $n+1$. Let us call this latent connection between the $n$-th domino and the $n+1$-th domino \textit{the inductive step}. But this alone doesn't really do anything. We need what we shall call \textit{the base case}: namely, domino 1 falling over. The base case and the inductive step together imply that all of the dominoes fall over. After all, the base case says that domino 1 falls over; and the $n=1$ instance of the inductive step connects domino 1 falling to domino 2 falling. But then we can repeatedly apply the inductive step: \textit{if} 2 falls over, then so does 3; and so on, and so on.

\textbf{Example 2}. The mathematician (and magician, and concert pianist, and philosopher, and...) Raymond Smullyan always gave people two pieces of advice: (1) always speak the truth; (2) each day, say `I will repeat this sentence tomorrow.' Why would he have done that?

Hopefully, you have some intuitive grip on these examples. Let's now try to formalize these ideas. Any proof by induction consists of two steps:
\begin{itemize}
    \item[] \textbf{The Base Case}. In this class, the base case will almost always involve proving something about the atomic formulas, for example, that they have a particular property. (Often, this step is basically trivial, so trivial that you might find it confusing at first.)
    \item[] \textbf{The Inductive Step}. In the second step, we pick some arbitrary formulas $A$ and $B$, and \textit{assume} that they have the desired property. (This is called the \textbf{inductive hypothesis}.) Then, we show that, given this assumption, \textit{any} formula we can build from $A$ and $B$ will inherit that property.\footnote{See the \textit{Induction Template} cheat sheet on Canvas.}
\end{itemize}
By these two steps, we're able to show that all formulas of propositional logic have the property of interest.

\textbf{3. But why does induction work?} Because our definition of well-formed formulas is recursive! Let's switch from thinking of a line of individual dominoes to thinking of a sequence of \textit{rows of dominoes}. We can think of these rows as being arranged in an inverted pyramid (see below).

Recall from Lecture 2:
\begin{itemize}
    \item[(i)] Every lower case letter (potentially with a subscript) is a formula of the language. For instance, $p$ is one and $r_5$ is another. \textbf{This is the bottom level of our inverted pyramid. All constructions begin here.}
    \item[(ii)] Let $A$ represent a generic expression of the language. If $A$ is a formula of the language, then so is $\Neg A$. \textbf{This is one way of moving a step up the pyramid.}
    \item[(iii)] Let $A$ and $B$ represent two generic expressions of the language. If $A$ and $B$ are formulas of the language, then:
    \begin{itemize}
        \item $(A \Conj B)$ and $(A \Disj B)$ and $(A \Cond B)$ and $(A \Bicond B)$ are each formulas of the language.
    \end{itemize}
    \textbf{These are all ways of using two formulas on the pyramid, $A$ and $B$, to build yet more formulas higher on the pyramid.}
    \item[(iv)] Nothing else is a formula of the language. \textbf{Crucially, this part of the definition says, that there are no other ways to build formulas. That is, every formula must have come from repeated applications of the above steps.}
\end{itemize}
Would induction work without clause (iv)? No! Without (iv) our base case would still tell us that the first rung of dominoes falls over. Moreover, our inductive step would still tell us that all of the dominoes \textit{connected to the first rung} will fall over. But are we sure that \textit{literally all of the dominoes} are linked to the first rung in this way? The purpose of clause (iv) is to guarantee this.

\begin{center}
\begin{tikzpicture}[x=1cm, y=1.2cm, font=\small]
    % Draw the outer truncated inverted pyramid (trapezoid)
    \draw[thick, fill=gray!5] (-4.5,4) -- (4.5,4) -- (1.5,0) -- (-1.5,0) -- cycle;

    % Horizontal tier dividers
    \draw[thick] (-3.75,3) -- (3.75,3);
    \draw[thick] (-3.0,2) -- (3.0,2);
    \draw[thick] (-2.25,1) -- (2.25,1);

    % TIER 4 (Top / Widest): Complex Formulas
    \node[align=center] at (0,3.5) {
        \textbf{More Complex Formulas} \\
        $\Neg A, \ (A \Conj B), \ (A \Disj B), \ (A \Cond B), \ (A \Bicond B)$
    };

    % TIER 3: Intermediate steps / Formulas A and B
    \node[align=center] at (0,2.5) {
        \textbf{Arbitrary Formulas} \\
        $A, \ B$
    };

    % TIER 2: Ellipses showing construction layers
    \node[align=center] at (0,1.5) {
        $\vdots$ \\
        (Construction Steps)
    };

    % TIER 1 (Bottom / Narrowest): Atomic Formulas
    \node[align=center] at (0,0.5) {
        \textbf{Atomic Formulas} \\
        $p, \ q, \ r, \ \dots$
    };
\end{tikzpicture}
\end{center}

\textbf{4. Some Concrete Examples of Mathematical Induction}

\textbf{Example 3}. Let $C$ be any formula of propositional logic. \textit{Claim:} Then the number of occurrences of atomic formulas in $C$ is one greater than the number of occurrences of binary connectives in $C$. (Recall: We are here counting the \textit{binary} connectives: $\ConjTight, \Disj, \Cond, \Bicond$ but not $\Neg$).

\begin{proof}
Let $\text{atom}(C)$ be the number of occurrences of atomic formulas in $C$, and let $\text{conn}(C)$ be the number of occurrences of binary connectives in $C$. We want to show that for any formula $C$, we have $\text{atom}(C) = \text{conn}(C) + 1$.

\textit{Base Case}. The base case is easy. Suppose $A$ is just an atomic formula, e.g., $p$ or $q$, etc. Note that $p$ contains just one atomic formula (itself) and zero binary connectives. So, $\text{atom}(p) = 1$ and $\text{conn}(p) = 0$. Clearly, the desired property holds, since $1 = 0 + 1$. This same reasoning holds for all atomic formulas. This establishes the base case. Metaphorically, this is the first domino falling over. What other dominoes does it knock over? We need the inductive step.

\textit{Inductive Step}. Now suppose that $A$ and $B$ are arbitrary formulas of propositional logic, and assume that the desired property holds for both of them. This is the inductive hypothesis which we get for free:
\begin{itemize}
    \item $\text{atom}(A) = \text{conn}(A) + 1$
    \item $\text{atom}(B) = \text{conn}(B) + 1$
\end{itemize}
Metaphorically, this is our assumption that the $n$-th rung of dominoes will fall over. Our goal is then to show that the next rung of dominoes ($n+1$) will also fall over. That is, what we want to show is that, given these assumptions about $A$ and $B$, the complex formulas, $\Neg A$, $(A \Conj B)$, $(A \Disj B)$, $(A \Cond B)$, and $(A \Bicond B)$ all have the desired property as well.

Let's start with $\Neg A$. Concretely, our goal is to show that
$$\text{atom}(\Neg A) = \text{conn}(\Neg A) + 1$$
To achieve this goal, note that the number of atomic formulas in $\Neg A$ is the same as in $A$, so $\text{atom}(\Neg A) = \text{atom}(A)$. For instance, consider the case when $A=((p\Disj q)\Cond p)$ such that $\Neg A=\Neg ((p\Disj q)\Cond p)$, both have three atomic formulas. Similarly, the number of binary connectives is unchanged, so $\text{conn}(\Neg A) = \text{conn}(A)$. In our example, both have two binary connectives. By the inductive hypothesis, $\text{atom}(A) = \text{conn}(A) + 1$, so it follows immediately that $\text{atom}(\Neg A) = \text{conn}(\Neg A) + 1$, as desired.

Now let's consider $(A \Conj B)$. Concretely, our goal is to show that
$$\text{atom}(A \Conj B)=\text{conn}(A \Conj B) +1$$
To show this, let us first note that the number of atomic formulas is the sum of those in its parts: $\text{atom}(A \Conj B) = \text{atom}(A) + \text{atom}(B)$, e.g, if $A$ had five and $B$ had three then $(A \Conj B)$ would have eight atomic formulas. Using our inductive hypothesis, we can substitute for $\text{atom}(A)$ and $\text{atom}(B)$:
\begin{align}
\nonumber
\text{atom}(A \Conj B) &= \text{atom}(A) + \text{atom}(B) & \text{our counting arg.}\\
\nonumber
&=(\text{conn}(A) + 1) + (\text{conn}(B) + 1) & \text{inductive hyp.}\\
\nonumber
&= \text{conn}(A) + \text{conn}(B) + 2    & \text{simplification}
\end{align}
Now let us compute the number of binary connectives in $(A \Conj B)$. It is just the sum of those in its parts, plus one more for the $\ConjTight$ in the middle. Concretely, if $A$ had four binary connectives and $B$ had two then $(A \Conj B)$ would have seven of them. In general, we have,
$$\text{conn}(A \Conj B) = \text{conn}(A) + \text{conn}(B) + 1$$
Comparing our last two calculations yields the desired result: $\text{atom}(A \Conj B)=\text{conn}(A \Conj B) +1$. Hence, the desired property holds for $(A \Conj B)$. The other three cases ($\Disj$, $\Cond$, and $\Bicond$) are proved in a way more-or-less identical to this. Metaphorically, we have just established that (having assumed that domino $n$ falls over) that domino $n+1$ will also fall over.

Importantly, the cascade is then automatic. We do not have to re-prove that $n+1$ knocks over $n+2$. Our argument that each domino knocks over the next began with a \textit{completely generic} domino. Hence, we only have to make the argument once and then apply it over and over. Concretely, because our initial choice of the formulas $A$ and $B$ was arbitrary the proof is complete.
\end{proof}
%\vspace{-0.3cm}

\textbf{Example 4}. Let $C$ be any formula of propositional logic which does not contain the connective $\Neg$. \textit{Claim:} All such formulas, $C$, are true in the all-true model. \textbf{Note that our goal in this case is not to knock down all of the dominoes, but only those which do not contain $\Neg$.}
\vspace{-0.3cm}\begin{proof}  \textit{Base Case}. The base case is easy. Suppose $A$ is just an atomic formula (e.g., $p$, $q$, etc) all of which don't contain $\Neg$. Let's look at the truth table for $p$.
\begin{table}[h!]
        \centering
        \begin{tabular}{c|c}
             $p$ & $p$ \\
             \hline
             T & T\\
             F & F
        \end{tabular}
    \end{table}

Clearly, $p$ is true in the all-true model. The same goes for every atomic formula (our choice of $p$ was arbitrary). This establishes the base case. Once again, this is us establishing that the first domino falls over. Now let's investigate its down-stream consequences.

\textit{Inductive Step}. Now suppose that $A$ and $B$ are arbitrary formulas of propositional logic \textit{which do not contain $\Neg$}. We get to assume (for free) that they have the desired property. Namely, $A$ is true in the all-true model. Similarly, for $B$.

Metaphorically, this is our assumption that \textit{part of} the $n$-th rung of dominoes will fall over. Our goal is then to show that \textit{part of} the next rung of dominoes ($n+1$) will also fall over. Concretely, we want to show that, given these assumptions about $A$ and $B$, the complex formulas $(A \Conj B)$, $(A \Disj B)$, $(A \Cond B)$, and $(A \Bicond B)$ all have the desired property as well. Note: We are leaving out $\Neg A$ here because it falls outside of the scope of formulas we are considering; it contains a $\Neg$.

Let's start with $(A \Conj B)$. By the inductive hypothesis, $A$ and $B$ are both true in the all-true model. When $A$ and $B$ are both true, so too is $(A \Conj B)$. Hence, $(A \Conj B)$ is also true in the all-true model.

The other three cases ($\Disj$, $\Cond$, and $\Bicond$) are proved very similarly.

Let's reflect. The base case showed that all of the atomic formulas have the desired property. The inductive step then extends this conclusion to everything which can be built from repeatedly applying the $\ConjTight$, $\Disj$, $\Cond$, and $\Bicond$ constructions. But does ``everything build-able without the negation rule'' cover our desired target of ``everything which lacks a $\Neg$''? Yes, because if a formula has a $\Neg$ it must have come from applying the negation rule.\footnote{This bit of reasoning is implicitly applying the contrapositive rule: $p\Cond q$ is logically equivalent to $\Neg q\Cond \Neg p$. See Lecture 4 (and 6).}
\end{proof}

\begin{framed}
\noindent\textbf{Deep Thought: Functional Completeness.} This result has an interesting consequence. As we will see in Problem Set 2, every truth table can be expressed in DNF form using only the symbols $\{\ConjTight,\Disj,\Neg\}$. Could we remove $\Neg$ from this set and still express everything? What if we added the two conditional symbols back in $\{\ConjTight,\Disj,\Cond,\Bicond\}$?

Example 4 shows that even this set is insufficient. Suppose that I give you a sequence of eight truth values (e.g., F, F, T, T, F, T, F, F) and task you with designing a formula over $p$, $q$, and $r$ with the corresponding truth table. Since the truth table that I want you to construct is false in its first row, the corresponding formula \textit{must} have a $\Neg$ in it. The fact that we included something like $\Neg$ (which is false in its first row) in our language is non-negotiable.

Problem Set 2 will explore the topic of \textbf{functional completeness}: what is the minimal expressive toolkit? Does $\{\Neg, \ConjTight\}$ alone suffice? What about $\{\Neg, \Disj\}$? Can any one of our five symbols do everything alone? If not, is there some other connective which can? Ultimately, we chose our five connectives for readability; we could have chosen fewer. The trade-off is conciseness versus intelligibility: a design choice, not a necessity.
\end{framed}

%\textbf{Challenge Question}. Prove the following claim using mathematical induction. Let $A$ be a formula of propositional logic. If $A$ is built only from $\Neg$ and $\Bicond$ then it is either a tautology (always true) or it is a contradiction (always false) or it is balanced (true in half of the cases and false in half of the cases).

%\textbf{Challenge Question}. Prove that no formula of propositional logic contains two atomic formulas in a row (e.g., $pp$ is not a formula of propositional logic).

\end{document}